\chapter{Electromagnetic Field}

\begin{remarkbox}
Electromagnetism is the prototype of gauge field theory. It contains all the
central ideas that later reappear in Yang--Mills theory: potentials, gauge
freedom, field strength, an action principle, coupling to conserved matter
currents, constraints, propagating transverse degrees of freedom, local energy
and momentum, waves, and radiation.
\end{remarkbox}

\section{Maxwell Equations in Vector Form}

In units where the electric and magnetic constants are absorbed into the
definitions of fields and sources, Maxwell's equations are
\[
    \nabla\cdot\vect{E}=\rho,
\]
\[
    \nabla\cdot\vect{B}=0,
\]
\[
    \nabla\times\vect{E}+\partial_t\vect{B}=0,
\]
and
\[
    \nabla\times\vect{B}-\partial_t\vect{E}=\vect{J}.
\]
Here \(\rho\) is the charge density and \(\vect{J}\) is the current density.
The first and fourth equations contain sources. The second and third are
homogeneous equations and express the absence of magnetic monopoles and the
Faraday induction law.

The equations imply charge conservation. Taking the divergence of
\[
    \nabla\times\vect{B}-\partial_t\vect{E}=\vect{J}
\]
gives
\[
    \nabla\cdot\vect{J}
    =\nabla\cdot(\nabla\times\vect{B})-\partial_t(\nabla\cdot\vect{E}).
\]
Since \(\nabla\cdot(\nabla\times\vect{B})=0\) and
\(\nabla\cdot\vect{E}=\rho\), one obtains
\[
    \partial_t\rho+\nabla\cdot\vect{J}=0.
\]
Thus Maxwell's equations are consistent only with locally conserved charge.

\begin{definitionbox}
The electromagnetic field is described by an electric field \(\vect{E}\) and a
magnetic field \(\vect{B}\) satisfying Maxwell's equations. In relativistic
field theory these fields are not independent objects but components of a
single antisymmetric spacetime tensor.
\end{definitionbox}

The vector form is physically transparent and useful for calculations in a
chosen inertial frame. Its limitation is that Lorentz covariance is hidden. The
covariant formulation reveals that electric and magnetic fields are frame-dependent
parts of a single electromagnetic field strength.

\section{Electromagnetic Potentials}

The homogeneous Maxwell equations can be solved by introducing potentials. The
equation
\[
    \nabla\cdot\vect{B}=0
\]
implies, at least locally in a topologically simple region, that
\[
    \vect{B}=\nabla\times\vect{A},
\]
where \(\vect{A}\) is the vector potential. Substituting this into Faraday's law
\[
    \nabla\times\vect{E}+\partial_t\vect{B}=0
\]
gives
\[
    \nabla\times\left(\vect{E}+\partial_t\vect{A}\right)=0.
\]
Therefore there exists a scalar potential \(\Phi\) such that
\[
    \vect{E}+\partial_t\vect{A}=-\nabla\Phi.
\]
Thus
\[
    \vect{E}=-\nabla\Phi-\partial_t\vect{A},
    \qquad
    \vect{B}=\nabla\times\vect{A}.
\]

The potentials are naturally combined into the four-potential
\[
    A^\mu=(\Phi,\vect{A}).
\]
With the mostly-minus metric convention,
\[
    A_\mu=(\Phi,-\vect{A}).
\]
The electromagnetic field strength tensor is then defined by
\[
    F_{\mu\nu}=\partial_\mu A_\nu-\partial_\nu A_\mu.
\]
This definition automatically incorporates the homogeneous Maxwell equations.

\begin{remarkbox}
The potentials are not merely calculational conveniences. In the action
formulation, the potential \(A_\mu\), not \(\vect{E}\) and \(\vect{B}\)
separately, is the fundamental field variable.
\end{remarkbox}

The price of using potentials is redundancy: many different potentials describe
the same electric and magnetic fields. This redundancy is gauge freedom.

\section{Gauge Freedom}

The electric and magnetic fields are unchanged under the transformation
\[
    \vect{A}\longrightarrow \vect{A}'=\vect{A}+\nabla\chi,
\]
\[
    \Phi\longrightarrow \Phi'=\Phi-\partial_t\chi,
\]
where \(\chi(t,\vect{x})\) is an arbitrary smooth function. Indeed,
\[
    \vect{B}'=\nabla\times\vect{A}'
    =\nabla\times\vect{A}+\nabla\times\nabla\chi
    =\vect{B},
\]
and
\[
\begin{aligned}
    \vect{E}'
    &=-\nabla\Phi'-\partial_t\vect{A}' \\
    &=-\nabla(\Phi-\partial_t\chi)-\partial_t(\vect{A}+\nabla\chi) \\
    &=-\nabla\Phi-\partial_t\vect{A}
    =\vect{E}.
\end{aligned}
\]
In covariant notation, this is
\[
    A_\mu\longrightarrow A'_\mu=A_\mu+\partial_\mu\chi.
\]
Then
\[
    F'_{\mu\nu}=\partial_\mu A'_\nu-\partial_\nu A'_\mu=F_{\mu\nu}.
\]

\begin{definitionbox}
A \emph{gauge transformation} of the electromagnetic potential is
\[
    A_\mu\longrightarrow A_\mu+\partial_\mu\chi.
\]
It changes the potential but leaves the field strength \(F_{\mu\nu}\) unchanged.
\end{definitionbox}

Gauge freedom means that the potential contains more components than physical
degrees of freedom. This is not a defect. It is the organizing principle of the
theory. The redundant description makes locality and Lorentz covariance
compatible with the existence of massless spin-one radiation.

\section{Maxwell Theory from an Action}

The free electromagnetic action is
\[
    S[A]=\int d^{d+1}x\,\mathcal{L}_{\mathrm{EM}},
\]
with
\[
    \mathcal{L}_{\mathrm{EM}}
    =-\frac{1}{4}F_{\mu\nu}F^{\mu\nu}.
\]
In four spacetime dimensions this is the standard Lorentz-invariant Maxwell
Lagrangian density. The factor \(-1/4\) gives the conventional normalization of
the equations of motion.

Let us vary the action. Since
\[
    \delta F_{\mu\nu}=\partial_\mu\delta A_\nu-\partial_\nu\delta A_\mu,
\]
one finds, using antisymmetry of \(F^{\mu\nu}\),
\[
    \delta\mathcal{L}_{\mathrm{EM}}
    =-\frac{1}{2}F^{\mu\nu}\delta F_{\mu\nu}
    =-F^{\mu\nu}\partial_\mu\delta A_\nu.
\]
After integration by parts,
\[
    \delta S
    =\int d^{d+1}x\,\partial_\mu F^{\mu\nu}\delta A_\nu
    +\text{boundary term}.
\]
For variations that make the boundary term vanish, the field equation is
\[
    \partial_\mu F^{\mu\nu}=0.
\]
This is Maxwell's equation in vacuum.

If an external conserved current \(J^\mu\) is present, one adds the interaction
term
\[
    \mathcal{L}_{\mathrm{int}}=-J_\mu A^\mu.
\]
The field equation becomes
\[
    \partial_\mu F^{\mu\nu}=J^\nu.
\]

\begin{theorembox}
The Maxwell action
\[
    S[A]=\int d^{d+1}x\left(-\frac{1}{4}F_{\mu\nu}F^{\mu\nu}-J_\mu A^\mu\right)
\]
gives the sourced Maxwell equation
\[
    \partial_\mu F^{\mu\nu}=J^\nu
\]
provided the external current is conserved and boundary terms are controlled.
\end{theorembox}

Gauge invariance of the source term requires current conservation. Under
\(A_\mu\to A_\mu+\partial_\mu\chi\),
\[
    \delta S_{\mathrm{int}}
    =-\int d^{d+1}x\,J^\mu\partial_\mu\chi
    =\int d^{d+1}x\,\chi\partial_\mu J^\mu
    +\text{boundary term}.
\]
Thus gauge invariance requires
\[
    \partial_\mu J^\mu=0.
\]

\section{Field Strength Tensor \texorpdfstring{\(F_{\mu\nu}\)}{Fmunu}}

The electromagnetic field strength is the antisymmetric tensor
\[
    F_{\mu\nu}=\partial_\mu A_\nu-\partial_\nu A_\mu.
\]
It has six independent components in four spacetime dimensions. These components
are the three electric and three magnetic components.

With the conventions
\[
    A^\mu=(\Phi,\vect{A}),
    \qquad
    A_\mu=(\Phi,-\vect{A}),
\]
one has
\[
    F_{0i}=E_i.
\]
The purely spatial components encode the magnetic field. A common convention is
\[
    F_{ij}=-\epsilon_{ijk}B_k,
\]
with signs depending on the precise placement of indices and the convention for
\(\epsilon_{ijk}\). What is invariant is the tensorial object \(F_{\mu\nu}\), not
a particular matrix convention.

The two Lorentz scalars built from the electromagnetic field are
\[
    F_{\mu\nu}F^{\mu\nu}=2(|\vect{B}|^2-|\vect{E}|^2),
\]
and
\[
    F_{\mu\nu}\widetilde{F}^{\mu\nu}=-4\vect{E}\cdot\vect{B},
\]
where the dual field strength is
\[
    \widetilde{F}^{\mu\nu}=\frac{1}{2}\epsilon^{\mu\nu\rho\sigma}F_{\rho\sigma}.
\]
These invariants show that \(\vect{E}\) and \(\vect{B}\) are frame-dependent
parts of one field, while certain combinations are Lorentz invariant.

\begin{definitionbox}
The electromagnetic field strength tensor \(F_{\mu\nu}\) is gauge invariant and
antisymmetric. It is the curvature associated with the gauge potential \(A_\mu\).
\end{definitionbox}

In differential-form notation, the potential is a one-form
\[
    A=A_\mu dx^\mu,
\]
and the field strength is the two-form
\[
    F=dA.
\]
Gauge transformations are
\[
    A\longrightarrow A+d\chi,
\]
and the identity \(d^2=0\) gives
\[
    dF=0.
\]

\section{Maxwell Equations in Covariant Form}

The inhomogeneous Maxwell equations are
\[
    \partial_\mu F^{\mu\nu}=J^\nu.
\]
The homogeneous Maxwell equations are encoded by the Bianchi identity
\[
    \partial_{[\lambda}F_{\mu\nu]}=0.
\]
Equivalently,
\[
    \partial_\lambda F_{\mu\nu}
    +\partial_\mu F_{\nu\lambda}
    +\partial_\nu F_{\lambda\mu}=0.
\]
In differential-form notation this is simply
\[
    dF=0.
\]

The covariant equations are therefore
\[
    dF=0,
    \qquad
    \partial_\mu F^{\mu\nu}=J^\nu.
\]
In four-dimensional differential-form notation, the second equation can be
written schematically as
\[
    d\star F=\star J,
\]
where \(\star\) is the Hodge dual.

Charge conservation follows immediately from the inhomogeneous equation:
\[
    \partial_\nu J^\nu=\partial_\nu\partial_\mu F^{\mu\nu}=0,
\]
because \(F^{\mu\nu}\) is antisymmetric while \(\partial_\nu\partial_\mu\) is
symmetric.

\begin{summarybox}
Maxwell's four vector equations are compactly expressed as one Bianchi identity
and one sourced field equation:
\[
    \partial_{[\lambda}F_{\mu\nu]}=0,
    \qquad
    \partial_\mu F^{\mu\nu}=J^\nu.
\]
The first follows from \(F=dA\); the second follows from the Maxwell action.
\end{summarybox}

\section{Coupling to Charged Matter}

Electromagnetism couples to charged matter through a conserved current. At the
level of the action, the coupling is
\[
    S_{\mathrm{int}}=-\int d^{d+1}x\,J_\mu A^\mu.
\]
Gauge invariance requires
\[
    \partial_\mu J^\mu=0.
\]
This statement is not optional: a nonconserved external current would make the
gauge coupling inconsistent.

For a charged complex scalar field, the current is not imposed externally but is
constructed from the matter field. Using the covariant derivative
\[
    D_\mu\phi=(\partial_\mu+iqA_\mu)\phi,
\]
the scalar electrodynamics Lagrangian is
\[
    \mathcal{L}
    =-\frac{1}{4}F_{\mu\nu}F^{\mu\nu}
     +(D_\mu\phi)^*D^\mu\phi-m^2\phi^*\phi.
\]
Varying with respect to \(A_\mu\) gives a Maxwell equation with a matter current.
The current is covariantly built from \(\phi\), \(\phi^*\), and their gauge
covariant derivatives.

\begin{remarkbox}
The electromagnetic field does not couple to arbitrary matter variables. It
couples to a conserved \(U(1)\) current. This is why the previous chapter on the
complex scalar field is the natural bridge into Maxwell theory.
\end{remarkbox}

When matter is dynamical, conservation of the current is a consequence of the
matter field equations and gauge symmetry. When matter is external, conservation
must be imposed as a consistency condition.

\section{Lorenz Gauge and Coulomb Gauge}

Gauge freedom allows one to impose convenient conditions on \(A_\mu\). The
Lorenz gauge condition is
\[
    \partial_\mu A^\mu=0.
\]
It is Lorentz covariant and simplifies the Maxwell equation. Since
\[
    \partial_\mu F^{\mu\nu}
    =\partial_\mu(\partial^\mu A^\nu-\partial^\nu A^\mu)
    =\Box A^\nu-\partial^\nu(\partial_\mu A^\mu),
\]
the Lorenz gauge reduces the sourced equation to
\[
    \Box A^\nu=J^\nu.
\]
In vacuum,
\[
    \Box A^\nu=0.
\]

The Coulomb gauge condition is
\[
    \nabla\cdot\vect{A}=0.
\]
It is not manifestly Lorentz covariant, but it is useful for separating
constraint-like and propagating parts of the electromagnetic field. In Coulomb
gauge, the scalar potential is determined by the instantaneous Poisson equation
\[
    -\nabla^2\Phi=\rho,
\]
while the transverse part of \(\vect{A}\) contains the radiative degrees of
freedom.

\begin{definitionbox}
The \emph{Lorenz gauge} is \(\partial_\mu A^\mu=0\). The \emph{Coulomb gauge} is
\(\nabla\cdot\vect{A}=0\). They are gauge choices, not physical restrictions on
\(\vect{E}\) and \(\vect{B}\).
\end{definitionbox}

A gauge condition does not remove the physical electromagnetic field. It removes
redundancy in the potential description. Residual gauge transformations may
remain even after a gauge condition is imposed.

\section{Electromagnetic Stress-Energy Tensor}

The electromagnetic stress-energy tensor is
\[
    T^{\mu\nu}
    =-F^{\mu\rho}F^\nu{}_{\rho}
     +\frac{1}{4}\eta^{\mu\nu}F_{\rho\sigma}F^{\rho\sigma}.
\]
With the mostly-minus convention, this tensor gives the standard positive energy
density
\[
    T^{00}=\frac{1}{2}(|\vect{E}|^2+|\vect{B}|^2).
\]
It is symmetric, gauge invariant, and conserved in vacuum:
\[
    \partial_\mu T^{\mu\nu}=0.
\]
In the presence of matter, the electromagnetic stress-energy tensor alone is
not conserved; energy and momentum are exchanged between field and matter.

The trace of the electromagnetic stress-energy tensor in four spacetime
dimensions is zero:
\[
    T^\mu{}_{\mu}=0.
\]
This reflects the classical scale invariance of source-free Maxwell theory in
four dimensions.

\begin{definitionbox}
The electromagnetic stress-energy tensor is the local density and flux tensor of
energy and momentum carried by the electromagnetic field. It is constructed from
\(F_{\mu\nu}\), not directly from the gauge potential \(A_\mu\), and is therefore
gauge invariant.
\end{definitionbox}

The stress-energy tensor is the object that couples to gravity. In general
relativity, electromagnetic fields curve spacetime through their stress-energy,
not merely through their energy density.

\section{Energy, Momentum, and Poynting Vector}

The electromagnetic energy density is
\[
    u=\frac{1}{2}(|\vect{E}|^2+|\vect{B}|^2).
\]
The energy flux is the Poynting vector
\[
    \vect{S}=\vect{E}\times\vect{B}.
\]
The local energy conservation law is Poynting's theorem:
\[
    \partial_t u+\nabla\cdot\vect{S}=-\vect{J}\cdot\vect{E}.
\]
The right-hand side is the rate at which electromagnetic energy is transferred
to charged matter.

In vacuum, \(\vect{J}=0\), so
\[
    \partial_t u+\nabla\cdot\vect{S}=0.
\]
Thus electromagnetic energy is locally conserved and flows with flux
\(\vect{S}\).

The electromagnetic momentum density is
\[
    \vect{g}=\vect{E}\times\vect{B}
\]
in the same units. Thus the Poynting vector is also the momentum density in
relativistic units. This is consistent with the fact that electromagnetic waves
carry both energy and momentum.

\begin{examplebox}
For a plane electromagnetic wave in vacuum, \(\vect{E}\), \(\vect{B}\), and the
propagation direction are mutually orthogonal. The energy flux points in the
direction of propagation:
\[
    \vect{S}=\vect{E}\times\vect{B}.
\]
\end{examplebox}

The conservation laws of electromagnetism are not separate from Maxwell's
equations. They follow from the field equations and the stress-energy tensor,
just as Noether's theorem predicts from spacetime translation invariance.

\section{Electromagnetic Waves}

In vacuum, Maxwell's equations imply wave equations for both \(\vect{E}\) and
\(\vect{B}\). Taking the curl of Faraday's law gives
\[
    \nabla\times(\nabla\times\vect{E})
    +\partial_t(\nabla\times\vect{B})=0.
\]
Using \(\nabla\cdot\vect{E}=0\), the vector identity
\[
    \nabla\times(\nabla\times\vect{E})
    =\nabla(\nabla\cdot\vect{E})-\nabla^2\vect{E}
\]
and \(\nabla\times\vect{B}=\partial_t\vect{E}\), one obtains
\[
    \partial_t^2\vect{E}-\nabla^2\vect{E}=0.
\]
Similarly,
\[
    \partial_t^2\vect{B}-\nabla^2\vect{B}=0.
\]
Thus electromagnetic disturbances propagate at the speed of light in these
units.

A plane wave solution has the form
\[
    \vect{E}(t,\vect{x})=\vect{E}_0 e^{-i\omega t+i\vect{k}\cdot\vect{x}},
\]
\[
    \vect{B}(t,\vect{x})=\vect{B}_0 e^{-i\omega t+i\vect{k}\cdot\vect{x}}.
\]
Maxwell's equations imply
\[
    \omega^2=|\vect{k}|^2,
\]
\[
    \vect{k}\cdot\vect{E}_0=0,
    \qquad
    \vect{k}\cdot\vect{B}_0=0,
\]
and
\[
    \vect{B}_0=\frac{1}{\omega}\vect{k}\times\vect{E}_0.
\]
Therefore electromagnetic waves are transverse.

\begin{definitionbox}
A free electromagnetic wave has two transverse physical polarizations. The
longitudinal and time-like components of the potential are removed by gauge
freedom and constraints.
\end{definitionbox}

The transversality of electromagnetic radiation is one of the clearest physical
signatures of gauge theory. The potential has four components, but the free
massless electromagnetic field propagates only two local physical degrees of
freedom.

\section{Retarded Potentials and Radiation}

In Lorenz gauge, the sourced Maxwell equation becomes
\[
    \Box A^\mu=J^\mu.
\]
This equation can be solved using a retarded Green's function:
\[
    A^\mu(x)=\int d^{d+1}y\,G_{\mathrm{ret}}(x-y)J^\mu(y).
\]
The retarded prescription enforces causality: the potential at \(x\) depends
only on sources in the past light cone of \(x\).

For localized time-dependent sources, the far field contains radiation. The
radiative part falls off with distance in such a way that the energy flux
through a large sphere remains finite. Physically, radiation is the part of the
electromagnetic field that carries energy and momentum away to infinity.

\begin{remarkbox}
Static fields store energy around sources, but radiation transports energy to
large distances. The distinction is tied to boundary behavior: radiation is
identified by its flux through surfaces at infinity.
\end{remarkbox}

A complete treatment of radiation requires more detailed analysis of retarded
solutions, multipole expansion, and asymptotic fields. For the present course,
the key point is structural: hyperbolic field equations plus retarded boundary
conditions give causal propagation from sources to fields.

\section{Physical Degrees of Freedom of the Electromagnetic Field}

The four-potential \(A_\mu\) has four components. However, not all four describe
physical degrees of freedom. Gauge transformations
\[
    A_\mu\longrightarrow A_\mu+\partial_\mu\chi
\]
identify potentials that describe the same electromagnetic field. In addition,
not all components of Maxwell's equations are evolution equations; Gauss's law
acts as a constraint on initial data.

In vacuum, after accounting for gauge freedom and constraints, the electromagnetic
field has two propagating degrees of freedom at each spatial momentum. These are
the two transverse polarizations of light.

This counting is easiest to see for a plane wave. The condition
\[
    \vect{k}\cdot\vect{E}_0=0
\]
forces the electric field to be perpendicular to the propagation direction. A
vector perpendicular to \(\vect{k}\) in three-dimensional space has two
independent components. The magnetic field is then determined by
\[
    \vect{B}_0=\frac{1}{\omega}\vect{k}\times\vect{E}_0.
\]
Thus the independent radiative data consist of two polarizations.

\begin{summarybox}
Electromagnetism is the fundamental Abelian gauge theory. The gauge potential
\(A_\mu\) is redundant, while the field strength
\(F_{\mu\nu}=\partial_\mu A_\nu-\partial_\nu A_\mu\) is gauge invariant. Maxwell's
equations follow from the action \(-\frac{1}{4}F_{\mu\nu}F^{\mu\nu}\), coupling to
matter requires a conserved current, and the free electromagnetic field
propagates two transverse physical polarizations.
\end{summarybox}
